% V033 易错辨析 | 保留 trapbox 环境
\begin{trapbox}

\begin{description}[style=nextline,leftmargin=2.8em,labelsep=0.5em,itemsep=0.55em]
 \item[\textbf{\textcolor{CTrap}{陷阱一：把有向面积当普通面积。}}]
 错写 $S=\frac12(x_1y_2-x_2y_1)$。方法一对照图中，右侧行列式为 $-1$，普通面积仍为 $1/2$。
 \textbf{修正：} 面积非负，必须取 $\frac12|x_1y_2-x_2y_1|$；交换 $A,B$ 只改变符号、不改变面积。

 \item[\textbf{\textcolor{CTrap}{陷阱二：记成同类坐标相乘。}}]
 错式 $\frac12|x_1x_2-y_1y_2|$。
 \textbf{最小反例：} $A=(1,0)$、$B=(0,1)$，正确面积 $1/2$，错式为 $0$。
 \textbf{修正：} 第一项取 $A$ 的横坐标 $\times$ $B$ 的纵坐标，第二项相反，交叉后做差。

 \item[\textbf{\textcolor{CTrap}{陷阱三：水平或竖直辅助点落在延长线，却仍将两面积相加。}}]
 $A=(4,3)$、$B=(2,-1)$；水平法得 $E=(-4/3,-1)$，
 $S_{\triangle OAB}=S_{\triangle EAB}-S_{\triangle EOB}=5$，不能相加。
 \textbf{修正：} 先确认 $E$ 是否在线段 $OA$ 内；在延长线上用面积差，统一公式最后取绝对值。

 \item[\textbf{\textcolor{CTrap}{陷阱四：相似比随意写，或者分母为零。}}]
 竖直法 $y_E=(x_2/x_1)y_1$ 只在 $x_1\ne0$ 时可写；
 水平法 $x_E=(y_2/y_1)x_1$ 只在 $y_1\ne0$ 时可写。
 \textbf{最小反例：} $A=(0,2)$、$B=(3,1)$，若硬除 $x_1=0$ 出错；直接用竖直底 $OA$ 得 $S=3$。
 \textbf{修正：} $OA$ 已水平或竖直时直接底乘高，不需要交点；若 $A=O$ 则面积为零。

 \item[\textbf{\textcolor{CTrap}{陷阱五：把平行四边形面积误当三角形面积。}}]
 $OADB$ 面积 $|x_1y_2-x_2y_1|$，而 $\triangle OAB$ 只有它的一半。
 \textbf{最小反例：} $A=(1,0), B=(0,1)$，平行四边形面积 $1$，三角形面积 $1/2$。

 \item[\textbf{\textcolor{CTrap}{陷阱六：三个点不含原点，直接套原点公式。}}]
 $P=(1,2),Q=(4,5),R=(2,7)$；若直接拿 $Q,R$ 求原点三角形是 $\frac12|4\cdot7-2\cdot5|=9$，但真正 $\triangle PQR$ 面积为 $6$。
 \textbf{修正：} 同时减去同一个顶点坐标，将其平移到原点。

 \item[\textbf{\textcolor{CTrap}{陷阱七：把方法一当成行列式公式的“原始证明”。}}]
 若学生尚不知 $[\vec u,\vec v]=u_xv_y-u_yv_x$ 为什么成立，直接说“叉积模等于面积”只是\emph{引用}未解释的性质。
 \textbf{修正：} 用竖直与水平辅助线分别推导来源，再把行列式作为最快计算工具。

 \item[\textbf{\textcolor{CTrap}{陷阱八：求面积前必须先判断共线。}}]
 这是多余步骤。\textbf{修正：} $S=\frac12|x_1y_2-x_2y_1|$ 对共线、非共线都适用；$S=0$ 即已自动判定共线。
\end{description}
\end{trapbox}
